\documentclass[11pt]{article}

\usepackage[margin=1in]{geometry}
\usepackage{amsmath}
\usepackage{amssymb}
\usepackage{mathtools}
\usepackage{graphicx}
\usepackage{hyperref}
\usepackage{listings}

\hypersetup{
    hidelinks,
    pdftitle={Geometric Motion Planning: Theory and Algorithms},
    pdfauthor={Travis Haran}
}

\setcounter{tocdepth}{2}
\setlength{\parindent}{0pt}
\setlength{\parskip}{0.55em}

\lstset{
    basicstyle=\ttfamily\small,
    columns=fullflexible,
    keepspaces=true,
    frame=single,
    showstringspaces=false
}

\title{\textbf{Geometric Motion Planning}\\[0.35em]\large Theory and Algorithms}
\author{Travis Haran}
\date{\today}

\begin{document}

\pagenumbering{roman}
\maketitle
\thispagestyle{empty}
\clearpage

\tableofcontents
\clearpage

\pagenumbering{arabic}
\setcounter{page}{1}

\section{Introduction}

This document contains the theoretical foundations behind the
Geometric Motion Planning project.

The goal of the project is to develop a motion-planning framework
from first principles, beginning with fundamental computational
geometry operations and gradually progressing toward classical
robot motion planning, configuration-space methods, sampling-based
planning, and humanoid footstep planning.

The theory presented here is developed alongside the implementation.
Each major algorithm includes its mathematical foundation, geometric
interpretation, computational complexity, and connection to the C++
implementation.

\section{Geometric Primitives}

The motion-planning system begins with a small collection of geometric
primitives. These structures provide the representation used by all later
geometry and planning algorithms.

\subsection{Point}

A two-dimensional point is represented as
\[
p=(x,y), \qquad x,y\in\mathbb{R}.
\]

In the implementation, coordinates are stored using double-precision
floating-point values. A point is represented in C++ as:

\begin{lstlisting}[language=C++]
struct Point {
    double x;
    double y;
};
\end{lstlisting}

\subsection{Line Segment}

A line segment is defined by two endpoints $A$ and $B$ and is denoted by
$\overline{AB}$. If
\[
A=(x_A,y_A), \qquad B=(x_B,y_B),
\]
then the segment contains every point
\[
P(t)=A+t(B-A), \qquad 0\leq t\leq 1.
\]

The C++ representation is:

\begin{lstlisting}[language=C++]
struct Segment {
    Point a;
    Point b;
};
\end{lstlisting}

\subsection{Polygon}

A polygon is represented as an ordered sequence of vertices
\[
\mathcal{P}=\{p_0,p_1,\ldots,p_{n-1}\}.
\]

Each consecutive pair of vertices forms an edge, and the final vertex
connects back to the first. Thus the edges are
\[
\left(p_i,p_{(i+1)\bmod n}\right), \qquad 0\leq i<n.
\]

This modular indexing automatically generates the closing edge
$p_{n-1}\rightarrow p_0$.

For this project, a valid polygon contains at least three vertices. Polygons
with fewer than three vertices are treated as invalid and therefore produce no
polygon edges.

Obstacles are assumed to be \emph{simple polygons} unless otherwise stated. A
simple polygon has no self-intersections between non-adjacent edges.
Self-intersecting polygons are outside the current geometric model.

\subsection{Path}

A robot path is represented using an ordered list of waypoints
\[
P_0,P_1,\ldots,P_k.
\]

The corresponding path segments are
\[
(P_0,P_1),(P_1,P_2),\ldots,(P_{k-1},P_k).
\]

Storing waypoints rather than explicitly storing every segment avoids
duplicating information and simplifies path modification.

\section{Orientation and the 2D Cross Product}

One of the most important predicates in computational geometry is the
orientation test. Given three points
\[
A=(x_A,y_A),\qquad B=(x_B,y_B),\qquad C=(x_C,y_C),
\]
we want to determine whether traveling from $A$ to $B$ and then toward $C$
represents a left turn, right turn, or no turn.

Define
\[
\overrightarrow{AB}=B-A, \qquad \overrightarrow{AC}=C-A.
\]
The two-dimensional cross product is
\[
\overrightarrow{AB}\times\overrightarrow{AC}
=(x_B-x_A)(y_C-y_A)-(y_B-y_A)(x_C-x_A).
\]

We therefore define
\[
\operatorname{orientation}(A,B,C)
=(x_B-x_A)(y_C-y_A)-(y_B-y_A)(x_C-x_A).
\]

The sign gives the geometric relationship:
\[
\operatorname{orientation}(A,B,C)
\begin{cases}
>0, & \text{counterclockwise / left turn},\\
<0, & \text{clockwise / right turn},\\
=0, & \text{collinear}.
\end{cases}
\]

\subsection{Geometric Interpretation}

The magnitude of the cross product is the area of the parallelogram formed
by $\overrightarrow{AB}$ and $\overrightarrow{AC}$. Therefore the area of
triangle $ABC$ is
\[
\operatorname{Area}(ABC)
=\frac{1}{2}\left|\operatorname{orientation}(A,B,C)\right|.
\]

If the orientation is zero, the triangle has zero area, meaning that the
three points are collinear.

\subsection{Floating-Point Considerations}

When coordinates are represented using floating-point values, exact
comparison with zero can be unreliable. Instead, a small tolerance
$\varepsilon$ is used:
\[
|x|<\varepsilon.
\]

In this project,
\[
\varepsilon=10^{-9}.
\]

Values sufficiently close to zero are therefore treated as zero when
checking collinearity.

\section{Segment Intersection}

Given two line segments $\overline{AB}$ and $\overline{CD}$, we want to
determine whether they intersect. This operation is fundamental to collision
detection, polygon validation, visibility testing, and motion planning.

\subsection{Proper Intersections}

A proper intersection occurs when two line segments cross through the
interiors of both segments.

Let the two segments be

\[
\overline{AB}
\qquad\text{and}\qquad
\overline{CD}.
\]

Define

\[
o_1 = \operatorname{orientation}(A,B,C),
\]

\[
o_2 = \operatorname{orientation}(A,B,D),
\]

\[
o_3 = \operatorname{orientation}(C,D,A),
\]

and

\[
o_4 = \operatorname{orientation}(C,D,B).
\]

The segments properly intersect when the endpoints of each segment lie on
opposite sides of the other segment:

\[
o_1 o_2 < 0
\]

and

\[
o_3 o_4 < 0.
\]

A proper intersection therefore excludes endpoint touching and collinear
overlap. These cases may still count as general segment intersections, but
they are geometrically different from a crossing through both segment
interiors.

This distinction is useful for visibility testing. A candidate visibility
edge should be rejected when it properly crosses an obstacle boundary, while
some forms of boundary touching must remain allowed.

\subsection{Degenerate Cases}

The proper-intersection test alone does not handle endpoint touching,
collinear overlap, or one endpoint lying on another segment. These cases
require a point-on-segment test.

A point $P$ lies on $\overline{AB}$ if
\[
\operatorname{orientation}(A,B,P)=0
\]
and
\begin{align*}
\min(x_A,x_B)&\leq x_P\leq\max(x_A,x_B),\\
\min(y_A,y_B)&\leq y_P\leq\max(y_A,y_B).
\end{align*}

In the current project convention, endpoint touching and collinear overlap
both count as intersection.

\subsection{Complexity}

The segment-intersection test performs a constant number of arithmetic
operations. Therefore,
\[
T(n)=O(1).
\]

Although simple, this predicate becomes one of the most frequently used
operations in later motion-planning algorithms.

\section{Point-in-Polygon Testing}

Given a polygon $\mathcal{P}$ and query point $q$, we want to determine
whether $q$ lies inside the polygon. The implementation in this project uses
the winding-number algorithm.

\subsection{Winding Number}

Imagine tracing the polygon boundary while observing how many times the
polygon winds around the query point. For a simple polygon,
\[
WN(q)=
\begin{cases}
0, & \text{if } q \text{ is outside},\\
+1 \text{ or } -1, & \text{if } q \text{ is inside}.
\end{cases}
\]

The sign depends on whether the polygon vertices are ordered
counterclockwise or clockwise. The implementation therefore uses
\[
WN(q)\neq 0
\]
as the inside condition.

\subsection{Upward Crossings}

For an edge from $A$ to $B$, an upward crossing occurs when
\[
y_A\leq y_q \qquad \text{and} \qquad y_B>y_q.
\]
If additionally
\[
\operatorname{orientation}(A,B,q)>0,
\]
the winding number is incremented:
\[
WN\leftarrow WN+1.
\]

\subsection{Downward Crossings}

A downward crossing occurs when
\[
y_B\leq y_q \qquad \text{and} \qquad y_A>y_q.
\]
If
\[
\operatorname{orientation}(A,B,q)<0,
\]
then
\[
WN\leftarrow WN-1.
\]

The asymmetric inequalities prevent polygon vertices lying exactly on the
horizontal query line from being counted twice.

\subsection{Boundary Convention}

For collision detection, this project treats points lying on the polygon
boundary as inside the obstacle. Therefore, before updating the winding
number, each edge is checked using the point-on-segment predicate.

\subsection{Complexity}

Each polygon edge is examined exactly once. For a polygon containing $n$
vertices,
\[
T(n)=O(n).
\]

\section{Collision Checking}

The geometric predicates developed so far can now be combined to answer
questions relevant to robot motion.

\subsection{Segment--Polygon Boundary Intersection}

Given a segment $s$ and polygon $\mathcal{P}$, the segment is tested against
every polygon edge. If
\[
E(\mathcal{P})=\{e_0,e_1,\ldots,e_{n-1}\},
\]
then we evaluate
\[
\operatorname{segmentsIntersect}(s,e_i)
\]
for every edge $e_i$.

The segment intersects the polygon boundary if
\[
\exists e_i\in E(\mathcal{P}) \text{ such that } s\cap e_i\neq\varnothing.
\]
For a polygon containing $n$ edges, this requires $O(n)$ segment-intersection
tests.

\subsection{Why Boundary Intersection Is Not Enough}

Consider a segment whose two endpoints lie completely inside a polygon. The
segment may never cross the polygon boundary. In that case,
\[
\text{boundary intersection}=\text{false},
\]
even though the segment lies entirely inside the obstacle.

This demonstrates an important distinction between \emph{boundary
intersection} and \emph{collision}.

\subsection{Collision-Free Segment}

A segment is considered collision free only if, for every obstacle,
\begin{enumerate}
    \item the segment does not intersect the obstacle boundary,
    \item the first endpoint is not inside the obstacle, and
    \item the second endpoint is not inside the obstacle.
\end{enumerate}

Symbolically,
\[
\operatorname{CollisionFree}(s)
=\neg\bigl(
\operatorname{BoundaryIntersection}(s)
\lor \operatorname{Inside}(s.a)
\lor \operatorname{Inside}(s.b)
\bigr).
\]

Touching an obstacle boundary is currently considered a collision.

\subsection{Collision-Free Path}

A path consisting of waypoints
\[
P_0,P_1,\ldots,P_k
\]
contains the segments
\[
(P_i,P_{i+1}), \qquad 0\leq i<k.
\]

The path is collision free if every one of these segments is collision free:
\[
\operatorname{CollisionFreePath}(\mathcal{P})
=\bigwedge_{i=0}^{k-1}
\operatorname{CollisionFree}\!\left(\overline{P_iP_{i+1}}\right).
\]

This creates the first complete geometric motion-validity test in the
project. Later planners will repeatedly call these collision predicates when
constructing visibility graphs, PRMs, RRTs, and humanoid footstep transitions.

\section{Visibility Graphs}

A visibility graph converts a continuous geometric planning problem into a
discrete graph-search problem.

For polygonal obstacles, the graph will eventually contain the start point,
goal point, and obstacle vertices as nodes. Two nodes are connected when they
are geometrically visible to one another.

\subsection{Geometric Visibility}

Two points $A$ and $B$ are considered visible when the open segment between
them does not pass through the interior of any obstacle.

For an obstacle $O$, this can be expressed as

\[
(A,B) \cap \operatorname{Interior}(O) = \varnothing.
\]

The use of the open segment $(A,B)$ is important. A visibility segment may
terminate exactly at an obstacle vertex without entering the obstacle
interior.

Visibility therefore differs from the Phase 1 collision convention. In
ordinary collision checking, touching an obstacle boundary counts as a
collision. In visibility testing, touching or traveling along an obstacle
boundary may be allowed as long as the segment does not enter the obstacle
interior.

\subsection{Proper Boundary Crossings}

A candidate visibility segment is invalid if it properly intersects an
obstacle edge. A proper intersection crosses through the interiors of both
segments rather than merely touching at an endpoint.

Proper-intersection testing detects ordinary crossings through an obstacle,
but it is not sufficient by itself. A candidate segment can enter an obstacle
through one of its vertices without producing a proper intersection.

\subsection{Parametric Representation of a Segment}

A point along the segment from $A$ to $B$ can be represented parametrically
as

\[
P(t) = A + t(B-A),
\qquad 0 \leq t \leq 1.
\]

The parameter $t$ describes the position along the segment:

\[
P(0)=A,
\]

\[
P(1)=B,
\]

and

\[
P\left(\frac{1}{2}\right)
=
\frac{A+B}{2}.
\]

If a point $V$ is known to lie on the nonzero segment $\overline{AB}$, its
parameter can be recovered from either coordinate. For example,

\[
t = \frac{V_x-A_x}{B_x-A_x}
\]

when the $x$ component is suitable for division.

The implementation uses the coordinate with the larger absolute change
between $A$ and $B$. This avoids division by zero for vertical or horizontal
segments and reduces division by very small direction components.

The helper that computes this parameter assumes that the point already lies
on a nonzero segment. It does not perform segment-membership validation.

\subsection{Visibility Near Obstacle Vertices}

Proper-intersection testing alone cannot detect every invalid visibility
segment. In particular, a segment may pass through an obstacle vertex and
immediately enter the obstacle interior.

Suppose an obstacle vertex $V$ lies on a candidate segment and has parameter
$t_V$. Points immediately before and after the vertex can be sampled using

\[
P(t_V-\delta)
\]

and

\[
P(t_V+\delta),
\]

where $\delta$ is a small positive parameter displacement.

If either valid neighboring sample lies strictly inside the obstacle, then the
candidate segment enters the obstacle interior and is therefore not visible.

The endpoint cases require special treatment. If

\[
t_V=0,
\]

then $V=A$, so no portion of the candidate segment exists before the vertex.
Only the sample after the vertex needs to be considered. Similarly, if

\[
t_V=1,
\]

then $V=B$, so only the sample before the vertex needs to be considered.

The sampled parameters are also clamped to the valid segment interval

\[
0 \leq t \leq 1.
\]

This ensures that visibility testing considers only the finite candidate
segment rather than points on the infinite line extending beyond its
endpoints.

\subsection{Visibility Graph Construction}

Given a start point $S$, a goal point $G$, and a collection of polygonal
obstacles, the visibility graph is represented as

\[
\mathcal{G} = (V,E).
\]

The node set contains the start point, goal point, and every obstacle vertex:

\[
V =
\{S,G\}
\cup
\bigcup_{O \in \mathcal{O}}
\operatorname{Vertices}(O).
\]

In the current implementation, the start point is stored as node $0$, the
goal point as node $1$, and obstacle vertices are appended afterward.

To construct the edge set, every unique pair of nodes $(v_i,v_j)$ with

\[
i < j
\]

is tested for geometric visibility. If the two nodes are visible, an
undirected edge is added between them.

Thus,

\[
E =
\left\{
(v_i,v_j)
\mid
i < j
\text{ and }
\operatorname{Visible}(v_i,v_j)
\right\}.
\]

Using $i<j$ ensures that each unordered pair is considered exactly once. It
also prevents self-edges such as $(v_i,v_i)$.

\subsection{Euclidean Edge Weights}

Each visibility edge is assigned a weight equal to the Euclidean distance
between its endpoints.

For points

\[
A=(x_A,y_A)
\qquad\text{and}\qquad
B=(x_B,y_B),
\]

the edge weight is

\[
w(A,B)
=
\sqrt{
(x_B-x_A)^2
+
(y_B-y_A)^2
}.
\]

These weights represent the geometric length of traveling directly between
two visible nodes.

The resulting weighted graph can later be searched using shortest-path
algorithms such as Dijkstra's algorithm or A*.

\section{Shortest Paths with Dijkstra's Algorithm}

Once the weighted visibility graph has been constructed, motion planning
becomes a shortest-path problem on a graph.

For each node $v$, Dijkstra's algorithm maintains a tentative distance

\[
d(v),
\]

representing the shortest currently known distance from the start node to
$v$.

Initially,

\[
d(S)=0,
\]

while every other node is assigned

\[
d(v)=\infty.
\]

The central operation in Dijkstra's algorithm is \emph{edge relaxation}.
For an edge from $u$ to $v$ with weight $w(u,v)$, a new candidate distance
to $v$ is

\[
d(u)+w(u,v).
\]

If

\[
d(u)+w(u,v)<d(v),
\]

then a shorter route to $v$ has been discovered, and we update

\[
d(v)=d(u)+w(u,v).
\]

The predecessor of $v$ is also recorded as $u$ so that the final shortest
path can later be reconstructed.

\subsection{Selecting the Next Node}

After relaxing the edges of the current node, Dijkstra's algorithm selects
the unvisited node with the smallest tentative distance.

Once a node is selected in this way, its shortest-path distance is considered
finalized.

This property relies on all graph edge weights being non-negative. In the
visibility graph, edge weights are Euclidean distances, so this condition is
automatically satisfied.

\subsection{Path Reconstruction}

Whenever relaxation improves the tentative distance to a node $v$, the
algorithm records the node $u$ from which that improvement came.

This node is called the \emph{predecessor} or \emph{parent} of $v$.

After the goal is reached, the shortest path is reconstructed by following
the parent links backward:

\[
G \rightarrow \operatorname{parent}(G)
\rightarrow \operatorname{parent}(\operatorname{parent}(G))
\rightarrow \cdots \rightarrow S.
\]

Because this produces the path in reverse order, the resulting sequence is
reversed to obtain the final path from start to goal.

\subsection{Undirected Visibility Edges}

The visibility graph is undirected: if node $u$ can see node $v$, then node
$v$ can also see node $u$.

For efficiency, each undirected edge is stored only once in the graph.

Therefore, when processing a node, Dijkstra's algorithm must check whether
the current node appears as either endpoint of an edge.
\section{Breadth-First Search}

Breadth-First Search (BFS) is a graph-search algorithm that explores nodes
level by level outward from a start node.

BFS is designed for unweighted graphs, or equivalently for graphs in which
every edge has the same cost.

Its key property is that the first path it finds to a node uses the minimum
possible number of edges from the start.

Therefore, BFS minimizes

\[
\text{hop count},
\]

rather than weighted path cost.

\subsection{Queue-Based Exploration}

BFS uses a first-in, first-out queue.

The start node is inserted into the queue first.

The algorithm repeatedly:

\begin{enumerate}
    \item removes the node at the front of the queue,
    \item examines all of its unvisited neighbors,
    \item marks newly discovered neighbors as visited,
    \item records the current node as their parent,
    \item inserts those neighbors at the back of the queue.
\end{enumerate}

Because nodes are processed in the same order in which they are discovered,
the search expands the graph one level at a time.

Conceptually, the traversal proceeds as

\[
\text{level }0
\rightarrow
\text{level }1
\rightarrow
\text{level }2
\rightarrow
\cdots
\]

where level $k$ contains nodes that can be reached from the start using
exactly $k$ edges.

\subsection{Visited Nodes}

A visited structure is used to prevent nodes from being inserted into the
queue multiple times.

In the current project, graph nodes are indexed by dense integer IDs,

\[
0,1,\ldots,|V|-1,
\]

so a boolean array can be used:

\[
\texttt{visited}[i].
\]

A node is marked visited when it is inserted into the queue.

This prevents two different already-expanded nodes from inserting the same
neighbor into the queue more than once.

\subsection{Shortest Path in an Unweighted Graph}

Suppose the goal node is first discovered at level $k$.

Every node in levels

\[
0,1,\ldots,k-1
\]

has already been explored before any node at level $k+1$ is considered.

Therefore, no path containing fewer than $k$ edges can exist.

Hence BFS returns a path with minimum hop count.

If all graph edges have identical cost $c$, minimizing hop count also
minimizes total path cost because

\[
\text{cost}(P)
=
c \cdot
\text{number of edges in }P.
\]

\subsection{BFS on a Weighted Visibility Graph}

The visibility graph used in this project is weighted.

Each edge weight is its Euclidean length:

\[
w(u,v)
=
\|u-v\|.
\]

BFS ignores these edge weights during search.

It therefore optimizes

\[
\text{number of visibility edges}
\]

rather than

\[
\text{total Euclidean path length}.
\]

This means a BFS path can use fewer graph edges while still being
geometrically longer than the path found by Dijkstra's algorithm or A*.

For example, BFS may prefer a direct one-edge route of length

\[
10
\]

over a four-edge route whose total length is

\[
4.
\]

BFS chooses the first path because

\[
1 < 4,
\]

while Dijkstra chooses the second because

\[
4 < 10.
\]

This difference makes BFS useful in the project as a comparison algorithm,
even though it is not the correct weighted shortest-path planner for the
visibility graph.

\subsection{Path Reconstruction}

Like Dijkstra's algorithm and A*, BFS stores a parent for each newly
discovered node.

If node $v$ is first discovered while expanding node $u$, then

\[
\text{parent}(v)=u.
\]

Once the goal is reached, parent links are followed backward:

\[
G
\rightarrow
\text{parent}(G)
\rightarrow
\cdots
\rightarrow
S.
\]

This produces the path in reverse order.

Reversing the sequence gives

\[
S
\rightarrow
\cdots
\rightarrow
G.
\]

\subsection{Reporting Geometric Path Length}

Although BFS ignores edge weights during search, the implementation still
reports the total geometric length of the final path.

If the BFS path contains edges

\[
e_1,e_2,\ldots,e_k,
\]

then the reported geometric distance is

\[
L_{\text{BFS}}
=
\sum_{i=1}^{k} w(e_i).
\]

This value is measured after the path has been selected.

It does not influence the BFS search order.

Reporting geometric path length allows BFS results to be compared directly
with Dijkstra and A* in the interactive planner.

\subsection{Undirected Visibility Edges}

Visibility is symmetric.

If node $u$ is visible from node $v$, then node $v$ is also visible from
node $u$.

The visibility graph therefore represents an undirected graph.

Each undirected edge is stored only once, but the BFS adjacency list inserts
both directions:

\[
u \rightarrow v
\]

and

\[
v \rightarrow u.
\]

This allows BFS to traverse the visibility graph correctly regardless of
which endpoint appears first in the stored edge.

\subsection{Expanded Nodes}

For consistency with Dijkstra and A*, this project counts a node as expanded
when it has been removed from the queue and its neighbors are about to be
examined.

The goal node is not counted as expanded if the search terminates immediately
when the goal is removed from the queue.

This shared definition allows expanded-node counts to be compared across all
three planners.

\subsection{Complexity}

Using an adjacency-list representation, every vertex is inserted into and
removed from the queue at most once.

Every undirected edge is examined a constant number of times.

Therefore, the time complexity is

\[
O(|V|+|E|).
\]

The visited array, parent array, adjacency list, and queue require

\[
O(|V|+|E|)
\]

space overall.

If only the additional search state beyond the graph representation is
counted, the BFS working memory is

\[
O(|V|).
\]

\subsection{Comparison with Dijkstra and A*}

The three current graph-search planners differ primarily in how they choose
which node to explore.

BFS processes nodes according to graph depth and minimizes

\[
\text{hop count}.
\]

Dijkstra selects the node with the smallest known weighted cost

\[
g(n)
\]

and minimizes total path weight.

A* selects nodes using

\[
f(n)=g(n)+h(n),
\]

where $h(n)$ estimates the remaining cost to the goal.

For the current Euclidean visibility graph:

\begin{itemize}
    \item BFS is not guaranteed to produce the geometrically shortest path.
    \item Dijkstra and A* should produce the same optimal geometric path cost.
    \item A* may expand fewer nodes than Dijkstra because of its heuristic.
\end{itemize}

This makes BFS, Dijkstra, and A* useful together for demonstrating the effect
of different graph-search objectives and search strategies.
\section{Shortest Paths with Dijkstra's Algorithm}

Once the weighted visibility graph has been constructed, motion planning
becomes a shortest-path problem on a graph.

For each node $v$, Dijkstra's algorithm maintains a tentative distance

\[
d(v),
\]

representing the shortest currently known distance from the start node to
$v$.

Initially,

\[
d(S)=0,
\]

while every other node is assigned

\[
d(v)=\infty.
\]

The central operation in Dijkstra's algorithm is \emph{edge relaxation}.
For an edge from $u$ to $v$ with weight $w(u,v)$, a new candidate distance
to $v$ is

\[
d(u)+w(u,v).
\]

If

\[
d(u)+w(u,v)<d(v),
\]

then a shorter route to $v$ has been discovered, and we update

\[
d(v)=d(u)+w(u,v).
\]

The predecessor of $v$ is also recorded as $u$ so that the final shortest
path can later be reconstructed.

\subsection{Selecting the Next Node}

After relaxing the edges of the current node, Dijkstra's algorithm selects
the unvisited node with the smallest tentative distance.

Once a node is selected in this way, its shortest-path distance is considered
finalized.

This property relies on all graph edge weights being non-negative. In the
visibility graph, edge weights are Euclidean distances, so this condition is
automatically satisfied.

\subsection{Path Reconstruction}

Whenever relaxation improves the tentative distance to a node $v$, the
algorithm records the node $u$ from which that improvement came.

This node is called the \emph{predecessor} or \emph{parent} of $v$.

After the goal is reached, the shortest path is reconstructed by following
the parent links backward:

\[
G \rightarrow \operatorname{parent}(G)
\rightarrow \operatorname{parent}(\operatorname{parent}(G))
\rightarrow \cdots \rightarrow S.
\]

Because this produces the path in reverse order, the resulting sequence is
reversed to obtain the final path from start to goal.

\subsection{Undirected Visibility Edges}

The visibility graph is undirected: if node $u$ can see node $v$, then node
$v$ can also see node $u$.

For efficiency, each undirected edge is stored only once in the graph.

Therefore, when processing a node, Dijkstra's algorithm must check whether
the current node appears as either endpoint of an edge.
\section{A* Search}

A* is a shortest-path search algorithm that extends Dijkstra's algorithm
by incorporating an estimate of the remaining distance to the goal.

For each node $n$, A* considers three quantities:

\[
g(n),
\]

the known cost of the shortest path discovered so far from the start node
to $n$,

\[
h(n),
\]

a heuristic estimate of the remaining cost from $n$ to the goal, and

\[
f(n)=g(n)+h(n),
\]

the estimated total cost of a path from the start to the goal that passes
through $n$.

While Dijkstra's algorithm selects the next node using only $g(n)$, A*
selects the node with the smallest value of $f(n)$.

This allows A* to use information about the location of the goal when
deciding which part of the graph to explore.

\subsection{Euclidean-Distance Heuristic}

The visibility graph is embedded in two-dimensional Euclidean space.
Therefore, a natural heuristic is the straight-line distance from a node
to the goal.

For a node

\[
n=(x_n,y_n)
\]

and goal

\[
G=(x_G,y_G),
\]

the heuristic is

\[
h(n)
=
\sqrt{
(x_G-x_n)^2
+
(y_G-y_n)^2
}.
\]

This gives the shortest possible geometric distance between the node and
the goal without considering obstacles.

\subsection{Admissibility}

For A* to guarantee an optimal path, the heuristic should not overestimate
the true remaining shortest-path cost.

A heuristic satisfying

\[
h(n)\leq h^*(n)
\]

is called \emph{admissible}, where $h^*(n)$ denotes the true shortest-path
cost from node $n$ to the goal.

For the visibility graph, every edge weight is a Euclidean distance.
Any path from $n$ to the goal therefore consists of one or more Euclidean
line segments.

By the triangle inequality,

\[
\|n-G\|
\leq
\|n-v_1\|
+
\|v_1-v_2\|
+
\cdots
+
\|v_k-G\|.
\]

Therefore, the straight-line Euclidean distance to the goal can never
exceed the length of a valid path through the visibility graph.

Hence,

\[
h(n)=\|n-G\|
\]

is an admissible heuristic.

\subsection{Consistency}

A stronger property than admissibility is \emph{consistency}.

A heuristic is consistent if, for every graph edge between nodes $u$ and
$v$,

\[
h(u)
\leq
w(u,v)+h(v).
\]

For the Euclidean visibility graph,

\[
h(u)=\|u-G\|,
\]

\[
w(u,v)=\|u-v\|,
\]

and

\[
h(v)=\|v-G\|.
\]

The triangle inequality gives

\[
\|u-G\|
\leq
\|u-v\|+\|v-G\|.
\]

Therefore, the Euclidean-distance heuristic is also consistent.

With a consistent heuristic, once a node is selected for expansion using
the smallest $f$-value, its best path cost does not need to be reconsidered
under the standard A* formulation used in this project.

\subsection{Relaxation in A*}

Although A* selects nodes using

\[
f(n)=g(n)+h(n),
\]

edge relaxation still operates on the actual path cost $g(n)$.

For an edge from $u$ to $v$, the candidate path cost is

\[
g_{\text{candidate}}
=
g(u)+w(u,v).
\]

If

\[
g_{\text{candidate}}<g(v),
\]

then

\[
g(v)=g_{\text{candidate}},
\]

and the parent of $v$ is updated to $u$.

The heuristic does not change the actual cost of a path. It only influences
which node A* chooses to explore next.

\subsection{Relationship to Dijkstra's Algorithm}

Dijkstra's algorithm can be viewed as a special case of A* in which

\[
h(n)=0
\]

for every node.

In that case,

\[
f(n)
=
g(n)+0
=
g(n),
\]

so A* selects nodes using exactly the same criterion as Dijkstra's
algorithm.

Therefore,

\[
\boxed{
\text{Dijkstra}
=
\text{A* with } h(n)=0
}
\]

This relationship makes the two algorithms particularly useful to compare
within the same visibility graph.

\end{document}
